\chapter{Integración continua y despliegue}
\label{cap:cicd}

\section{El pipeline}

Un solo flujo, \pkpath{.github/workflows/ci.yml}, disparado en \pkcode{push} y en
\pkcode{pull\_request}, con \pkcode{concurrency} por rama y
\pkcode{cancel-in-progress}: un empujón nuevo cancela la ejecución anterior de la
misma rama. Dos jobs independientes que corren en paralelo.

\begin{figure}[htbp]
  \centering
  \includegraphics[width=\linewidth,height=0.78\textheight,keepaspectratio]{ci-pipeline.pdf}
  \caption{Los dos jobs y sus pasos, en orden. Generado desde
    \pkpath{diagramas/src/ci-pipeline.dot}.}
  \label{fig:ci}
\end{figure}

\subsection{Job \emph{API}}

Levanta un \emph{service container} \pkcode{postgres:18-alpine} con
\pkcode{pg\_isready} como healthcheck, y después ejecuta una secuencia que
reproduce el arranque real del sistema. Tres pasos existen por razones que no son
obvias y están comentadas en el propio flujo:

\begin{itemize}[leftmargin=1.1em,itemsep=3pt,topsep=3pt]
  \item \textbf{Crear los esquemas a mano.} Un service container no monta
        \pkpath{docker-entrypoint-initdb.d}, así que el paso ejecuta
        \pkpath{db/init/01-schemas.sql} —los dos esquemas y la extensión
        \pkcode{pg\_trgm}— tal como lo haría el fichero de compose.
  \item \textbf{Migrar el esquema \pkcode{auth} antes que el de dominio.} Las
        tablas de dominio tienen claves foráneas hacia \pkcode{auth."user"}; sin
        este paso la primera migración de Alembic no puede correr.
  \item \textbf{Sembrar un set del catálogo.} Las pruebas de servicio buscan
        cartas reales en vez de inventarlas, así que el catálogo tiene que estar.
        Un set basta y tarda unos ocho segundos.
\end{itemize}

Después, \pkcode{pytest -q}, \pkcode{mypy src} y \pkcode{ruff check .}.

\begin{pkdecision}[CI corre sin ninguna clave de modelo]
\pkcode{GOOGLE\_API\_KEY} se deja explícitamente vacía, con este comentario en el
flujo: «el agente es opcional, el chat responde 503 sin él, y un repositorio
público no debe necesitar una clave para ponerse en verde». Es la prueba
ejecutable de \pkid{RNF-07}.
\end{pkdecision}

\subsection{Job \emph{Web}}

Tres pasos: \pkcode{checkout}, \pkcode{npm ci} y
\pkcode{npx turbo run typecheck lint test --filter=web...}. El filtro acota la
ejecución al workspace de la web y a sus dependencias; las tareas de raíz
cubrirían también la API, cuyo utillaje vive en el otro job.

\section{Imágenes}

\par\addvspace{8pt}\noindent\begin{pkbox}{colspec={X[1,l]X[1.4,l]X[3,l]}}
  Servicio & Base & Notas \\
  api & \pkcode{uv:}\allowbreak\pkcode{python3.13-}\allowbreak\pkcode{bookworm-slim} & Capa de dependencias sólo con el \emph{lockfile} (\pkcode{uv sync --frozen --no-install-project --no-dev}), luego el código. \pkcode{UV\_COMPILE\_BYTECODE=1}. \\
  web & \pkcode{node:22-slim}, dos etapas & Etapa de construcción con \pkcode{npm ci} sobre los \pkcode{package.json} antes de copiar el código; etapa final que ensambla la traza \pkcode{standalone}. \\
\end{pkbox}\par\addvspace{8pt}

El arranque de cada imagen levanta el esquema que le pertenece. La web ejecuta
\pkcode{node apps/web/scripts/migrate-auth.mjs \&\& node apps/web/server.js}. La
API espera:

\begin{lstlisting}[language=bash]
# The domain tables carry foreign keys into `auth.user`, which Better Auth
# creates from the web service. Both start at once on a fresh deployment, so
# this waits for that schema rather than crash-looping.
for i in $(seq 1 30); do alembic upgrade head && break; \
  echo 'waiting for the auth schema'; sleep 5; done; \
  alembic current | grep -q . && exec fastapi run src/pokedex/main.py \
  --host 0.0.0.0 --port 8000
\end{lstlisting}

\begin{pkwarn}[La configuración de migración está duplicada a propósito]
\pkpath{apps/web/scripts/migrate-auth.mjs} repite la configuración de
\pkpath{lib/auth.ts} porque no puede importarla: la construcción
\pkcode{standalone} entrega \emph{bundles}, no el TypeScript fuente. El propio
archivo lo advierte: «un plugin añadido allí y olvidado aquí es una tabla que
esto nunca crea». Cualquier cambio en la configuración de Better Auth hay que
hacerlo en los dos sitios.
\end{pkwarn}

\section{Forma del despliegue}

Producción corre en \textbf{Railway}, con dos servicios y una base de datos:

\par\addvspace{8pt}\noindent\begin{pkbox}{colspec={X[1,l]X[3.4,l]}}
  Servicio & URL pública \\
  web & \pkpath{https://web-production-78a4b.up.railway.app} \\
  api & \pkpath{https://api-production-81ae.up.railway.app} \\
  db & PostgreSQL 18 gestionado, sólo accesible desde la red privada \\
\end{pkbox}\par\addvspace{8pt}

No hay configuración de proveedor versionada en el repositorio: ni
\pkpath{railway.json}, ni \pkpath{nixpacks.toml}, ni \pkpath{Procfile}. Lo que
está definido son las dos imágenes y el contrato de variables, y eso basta: el
sistema se despliega igual en cualquier anfitrión que ejecute contenedores
—Railway, Fly, Render o un Kubernetes— con la misma topología:

\begin{center}
\begin{tikzpicture}[font=\footnotesize,
  box/.style={draw=pkRule, fill=pkRaised, rounded corners=3pt, minimum height=1.15cm,
              minimum width=3.5cm, align=center, inner sep=6pt},
  arrow/.style={-{Stealth[length=4pt]}, pkMuted, line width=0.6pt}]
  \node[box, draw=pkRed, fill=pkRed!8] (web) {\pksb servicio \texttt{web}\\[2pt]
    \scriptsize imagen \pkpath{apps/web/Dockerfile}};
  \node[box, draw=pkRed, fill=pkRed!8, right=2.0cm of web] (api) {\pksb servicio \texttt{api}\\[2pt]
    \scriptsize imagen \pkpath{apps/api/Dockerfile}};
  \node[box, draw=tWater, fill=tWater!8, below=1.1cm of $(web)!0.5!(api)$] (db)
    {\pksb PostgreSQL 18\\[2pt]\scriptsize esquemas \texttt{auth} + \texttt{pokedex}};
  \draw[arrow] (web.south) -- node[left,pos=0.5,font=\scriptsize,color=pkMuted]
    {\ \ migra \texttt{auth}} (db.west|-web.south) -- (db.north west);
  \draw[arrow] (api.south) -- node[right,pos=0.5,font=\scriptsize,color=pkMuted]
    {migra \texttt{pokedex}\ \ } (db.east|-api.south) -- (db.north east);
  \draw[arrow] (api) -- node[above,font=\scriptsize,color=pkMuted]{JWKS} (web);
\end{tikzpicture}
\end{center}

\par\addvspace{8pt}\noindent\begin{pkbox}{colspec={X[0.6,l]X[4,l]}}
  Servicio & Variables que el anfitrión debe fijar \\
  api & \pkcode{DATABASE\_URL} (forma asyncpg), \pkcode{AUTH\_JWKS\_URL}, \pkcode{AUTH\_ISSUER}, \pkcode{AUTH\_AUDIENCE}, \pkcode{CORS\_ORIGINS}, \pkcode{PORT} \\
  web & \pkcode{DATABASE\_URL} (forma \pkcode{postgresql://}, es \pkcode{node-pg}), \pkcode{BETTER\_AUTH\_URL}, \pkcode{BETTER\_AUTH\_SECRET}, \pkcode{NEXT\_PUBLIC\_API\_URL}, \pkcode{PORT}, \pkcode{HOSTNAME=0.0.0.0} \\
\end{pkbox}\par\addvspace{8pt}

Cuatro puntos que deciden si el despliegue funciona a la primera:

\begin{enumerate}[leftmargin=1.4em,itemsep=4pt,topsep=3pt]
  \item \textbf{\pkcode{NEXT\_PUBLIC\_API\_URL} es un argumento de construcción.}
        Next lo hornea en el bundle del cliente, así que tiene que ser la URL
        \emph{pública} de la API, la que resuelve un navegador —no la interna del
        contenedor. Cambiarla obliga a reconstruir la imagen. Si el anfitrión
        sólo permite variables de ejecución, hay que pasarla como
        \emph{build arg} del contenedor.
  \item \textbf{El emisor y el JWKS son URLs distintas.} \pkcode{AUTH\_ISSUER} es
        lo que el token declara y lo que ve el navegador; \pkcode{AUTH\_JWKS\_URL}
        es por dónde el contenedor de la API alcanza al de la web. En una red
        interna serán diferentes, y eso es correcto.
  \item \textbf{Los dos DSN tienen formas distintas.} La API usa
        \pkcode{postgresql+asyncpg://}; la web usa \pkcode{postgresql://} porque
        su cliente es \pkcode{node-pg}.
  \item \textbf{\pkcode{CORS\_ORIGINS} tiene que incluir el origen público de la
        web}, porque el navegador llama a la API cruzando origen.
\end{enumerate}

Para poblar un despliegue nuevo:

\begin{lstlisting}[language=bash]
pokedex-sync base1 base2 base3     # catalogo
pokedex-seed --count 1000          # un mercado de coleccionistas
pokedex-demo                       # tres cuentas de demostracion
\end{lstlisting}
